\chapter{1934:Harold C. Urey}

\label{chap:chemistry1934}

The Nobel Prize in Chemistry for the year 1934 was awarded to Harold C. Urey.

\textbf{Motivation:}

for his discovery of heavy hydrogen

\section{Harold C. Urey}

\subsection*{Early Life and Education}

Harold Clayton Urey was born on 1893-04-29 in Walkerton, Indiana, United States.
He was an American physical chemist awarded the Nobel Prize in Chemistry in 1934 for his discovery of heavy hydrogen.

\subsection*{Career and Major Research}

Urey studied at the University of Montana and received his doctorate at the University of California, Berkeley, in 1923 under Gilbert N. Lewis.
He spent 1923-1924 at the Niels Bohr Institute in Copenhagen, then taught at the Johns Hopkins University and, from 1929, at Columbia University.
In late 1931 he and his colleagues Ferdinand G. Brickwedde and George M. Murphy discovered the isotope of hydrogen of mass 2, which Urey named deuterium, by distilling liquid hydrogen and identifying the heavy isotope spectroscopically.
During the Second World War he directed work at Columbia on the separation of uranium isotopes for the Manhattan Project.
After the war he was professor at the University of Chicago from 1945 to 1958 and then at the University of California, San Diego, until 1981.
He applied isotope chemistry to geology, geochemistry, and cosmochemistry, studied the origin of the Moon and of the solar system, and examined lunar samples brought back by the Apollo missions.

\subsection*{Nobel Prize-winning Work}

The work for which Harold C. Urey was awarded the Nobel Prize in Chemistry in 1934 was described by the Nobel Committee as: "for his discovery of heavy hydrogen".

Urey discovered heavy hydrogen in 1932, when he showed that ordinary hydrogen contains a stable isotope of mass 2 present in small proportion.
He gave this isotope the name deuterium, prepared samples of heavy water, and was awarded the prize for this discovery.

\subsection*{Significance and Impact}

The discovery of deuterium revealed that the isotopes of the lightest element can be separated and studied, and it opened the fields of isotope separation and isotope chemistry.
Heavy water became important as a moderator for nuclear reactors, and isotopes prepared by the methods Urey developed are now used throughout science and medicine.

\subsection*{Later Career and Honors}

Urey received the Davy Medal of the Royal Society in 1940, the National Medal of Science in 1964, and the Priestley Medal of the American Chemical Society in 1973, among many other honours.
The lunar crater Urey and asteroid 4716 Urey are named in his honour.
He died on 1981-01-05 in La Jolla, California, United States.

\subsection*{Legacy}

The legacy of Harold C. Urey endures through the lasting impact of his scientific contributions.
Isotope chemistry and isotope separation, fields he founded, continue to underpin nuclear science, geochemistry, and medicine.

\subsection*{Modern Developments}

Urey's discovery of deuterium gave chemistry a new isotope essential to nuclear energy, tracer studies, and modern physics. Deuterium is used to produce heavy water for nuclear reactors and is a key to understanding metabolic and geochemical processes, while deuterated compounds are standard tools in NMR spectroscopy and drug development. Urey's later work on isotope separation also made carbon and oxygen isotope geochemistry possible.

\subsection*{Selected Publications}

\begin{itemize}

\item Urey, H. C. (1931). "The Isotopes of Hydrogen." Physical Review.

\item Urey, H. C., Brickwedde, F. G., \& Murphy, G. M. (1932). "A Hydrogen Isotope of Mass 2 and Its Concentration." Physical Review.

\item Urey, H. C. (1952). "The Planets: Their Origin and Development." New Haven: Yale University Press.

\end{itemize}

\clearpage

\section*{Chapter Summary}

The 1934 prize recognized Harold C. Urey for his discovery of heavy hydrogen. The discovery of deuterium opened a new chapter in isotope chemistry and nuclear physics.
